\chapter{Kernel APDU Status Words}
\label{app:kernel-status-words}

This appendix inventories the status words generated by the current oXiDe SE
kernel, including its optional built-in diagnostic modules. It describes the
implemented behavior, not the complete status-word vocabulary of ISO/IEC 7816
or GlobalPlatform. Each value below is written in wire order: SW1 followed by
SW2, in hexadecimal. A suffix \texttt{xx} denotes a variable second byte.

\section{Scope and Interpretation}

The APDU transport, command dispatcher, registry management, secure-channel
engines and isolated execution boundary can all determine the final status.
Consequently, a status must be interpreted together with the command, its
phase and the configured kernel modules. In particular, \texttt{6985} is a
shared rejection status; it does not uniquely identify a fault or a shortage
of resources.

A Rustlet handler may return its own SW1/SW2, which the kernel forwards through
the APDU completion path. Rustlet-backed Security Domains may also supply
statuses through their dispatch interfaces. These application-defined values
are not a finite kernel-generated catalogue and are not reserved by this
appendix. Receiving a listed value therefore does not, by itself, prove that
the kernel generated it. A kernel panic, power loss or transport interruption
may prevent any complete response; absence of a response is not a status word.

\section{GlobalPlatform Boundary and Diagnostic Policy}

GP 2.3.1 defines common errors in section 11.1.3 and command-specific errors
in Chapter 11. INSTALL explicitly includes \texttt{6A84}; SELECT does not.
The common unspecified execution error is \texttt{6400}. These tables do not
define our private \texttt{6F01}, \texttt{6F02} or \texttt{6F03} diagnoses
\cite{gp-card}.

The implementation therefore keeps hardware faults, allocation failures and
flash-probe failures distinct internally, and uses private diagnoses on
application and diagnostic commands. The GP dispatcher normalizes runtime
failures to the applicable GP categories before secure-response generation.
This conversion does not affect the private flash-probe command filter.
Application panic and watchdog termination remain separate from hardware
faults, although they share the generic runtime status.

A one-to-one mapping between every internal cause and a GP response is not a
claim of this implementation. Standard statuses deliberately group causes.
Registry mutation APIs retain typed causes for missing references, authority
refusal, invalid state or data, exhausted capacity and persistence failure.
Publication failures are propagated to the response rather than acknowledged
as successful RAM-only changes. Single-object replacements restore the old RAM
record on failure; deletion publishes a filtered BOSS before invalidating live
RAM objects. A Rustlet Security Domain's explicit error status is retained by
the proxy, then normalized at the GP boundary before response protection.

INSTALL, LOAD and PUT KEY use \texttt{6A84} for capacity exhaustion and
\texttt{6581} for memory failure. DELETE includes \texttt{6581}, but not
\texttt{6A84}. STORE DATA includes \texttt{6A84}, but not \texttt{6581};
SET STATUS includes neither. Where the command table lacks the specific
memory category, the implementation uses the general execution error
\texttt{6400}, keeping the original cause distinct internally
(GP tables 11-10, 11-26, 11-55, 11-60, 11-78, 11-87 and 11-90).
This catalogue is not a declaration of complete GP conformance.

\section{Command, Registry and Secure-Channel Statuses}

\begin{description}
\item[\texttt{9000} --- Success.]
The command completed successfully. For a deferred response, this is returned
when the final response chunk is drained, provided that the saved completion
status was successful.

\item[\texttt{6300} --- Authentication failed.]
The authentication value supplied during secure-channel establishment was
not accepted; this includes the SCP03 host cryptogram check.

\item[\texttt{6310} --- More registry data available.]
A \texttt{GET STATUS} response contains only part of the requested registry
listing. Continue with the command's next-occurrence mechanism. This is
command-level pagination, distinct from transport-level \texttt{61xx}.

\item[\texttt{6581} --- Persistent memory failure.]
The registry could not initialize, erase, program, decode or publish persistent
storage. INSTALL, LOAD, DELETE and PUT KEY expose this GP memory-failure
category. STORE DATA and SET STATUS translate it to \texttt{6400}.
The private flash-probe diagnoses \texttt{6F02}/\texttt{6F03} remain unchanged.

\item[\texttt{6600} --- Certificate verification failed.]
An SCP11 certificate did not pass verification. An absent referenced key or
certificate is instead reported through the referenced-data path below.

\item[\texttt{6700} --- Wrong length.]
The command data length, encoded field length, or response/state buffer
requirement is unacceptable. This also covers bounded ABI-buffer staging
failures; it is not the status for failure to allocate a Rustlet's RAM.

\item[\texttt{6982} --- Security status not satisfied.]
A required security proof is missing or invalid, or secure-message verification
rejects the command. This includes protected-command integrity failures,
insufficient management privileges and attempts outside the authority's subtree.

\item[\texttt{6985} --- Conditions not satisfied.]
The current state does not permit the operation. Uses include unmet lifecycle
conditions, dependency restrictions, unavailable execution or isolation
configurations, invalid startup descriptors, and secure-channel sequence or
session-state errors. Registry capacity, persistence and authority failures
are encoded separately as described above. Startup failure
before a valid Rustlet descriptor is obtained remains in this category,
while recovered hardware faults use the execution-error path.
Insufficient memory on \texttt{INSTALL} is reported separately as \texttt{6A84}.
For \texttt{SELECT}, unavailable activation retains \texttt{6985}.

\item[\texttt{6A80} --- Wrong data.]
The command payload is malformed or its contents are unacceptable. Examples
include invalid management encodings, inconsistent load blocks, and an
incompatible or invalid FAE rejected during loading.

\item[\texttt{6A82} --- File or application not found.]
The requested application, executable reference or AID cannot be resolved.
An unknown AID on \texttt{SELECT} is a typical case.

\item[\texttt{6A84} --- Insufficient memory or registry capacity.]
The registry has no free object slot, persistent storage has no safe reusable
span, or Rustlet activation cannot obtain its required RAM. This covers the writable-data budget,
the aligned stack and data allocation, and the kernel-side heap metadata. The allocation may fail because no suitably
sized and aligned free block remains, even when the sum of free bytes seems
sufficient. Activation produces this diagnosis, including reactivation of a displaced
instance. The GP dispatcher preserves it for \texttt{INSTALL}, but maps it to
\texttt{6985} for \texttt{SELECT}.
It identifies exhausted resources, not a fault in a running handler or a
failed flash write.

\item[\texttt{6A86} --- Incorrect P1/P2.]
The command does not accept the supplied parameter combination.

\item[\texttt{6A88} --- Referenced data not found.]
The requested registry record, data object or key reference is absent.
In particular, \texttt{GET STATUS} uses this when no matching record is found.

\item[\texttt{6D00} --- Instruction not supported.]
No enabled handler supports the instruction or management operation. The
T=0 front end also rejects INS values in the \texttt{6x} and \texttt{9x}
classes, which would be ambiguous with status-byte classes.

\item[\texttt{6400} --- GP execution failure.]
At the GP command-dispatch boundary, runtime statuses \texttt{6F00} and
\texttt{6F01} are mapped to this general GP category. Their internal causes
remain distinct, but the GP response does not encode that distinction.
This is also the fallback for a registry memory error whose specific SW is
absent from the current command's GP error table.

\item[\texttt{6F00} --- Internal execution failure.]
Runtime failures on application commands use this status. This includes
explicit Rustlet panic and watchdog termination. Built-in cryptographic diagnostics also
use it for internal failures. Recovered hardware isolation faults use
\texttt{6F01} instead.

\item[\texttt{6F01} --- Isolated hardware execution fault.]
A recoverable hardware trap terminated an isolated Rustlet: for example an
invalid memory access, an invalid instruction or a protected stack overflow.
This diagnosis covers both startup and handler execution. It is exposed on
application commands; the GP dispatcher maps it to \texttt{6400}.
It is an oXiDe SE private diagnostic, not a GP-defined status.
\end{description}

\section{Transport Completion Statuses}

\begin{description}
\item[\texttt{61xx} --- Response bytes available.]
Use \texttt{GET RESPONSE} to retrieve the buffered response.
SW2 gives the remaining byte count; \texttt{00} encodes 256 bytes in the
current short-APDU buffer model. The transport retains the command's completion
status and returns it after the last chunk, which need not be \texttt{9000}.

\item[\texttt{6Cxx} --- Correct response length.]
For a direct outgoing response, the requested length does not match the
produced length. SW2 gives the required length, with \texttt{00} encoding
256 bytes. This is a length-correction response, not a queued
\texttt{GET RESPONSE} continuation.
\end{description}

The transport also generates \texttt{6985} when \texttt{GET RESPONSE} has no
pending response, or another command arrives while response bytes remain
pending. These completion rules can replace or defer the status initially
returned by a command handler.

\section{Optional Diagnostic Module Statuses}

The \texttt{flash-probe} kernel module adds two private statuses to its
\texttt{INS=8E} diagnostic command family:
\begin{description}
\item[\texttt{6F02} --- Flash erase failed.]
The flash probe could not erase its reserved test sector.
\item[\texttt{6F03} --- Flash write failed.]
The flash probe could not program its test page.
\end{description}
These values are module-specific diagnostics, not general registry persistence
statuses. The same module uses \texttt{6985} when no usable probe area is
available. Other built-in test and monitoring modules reuse the common
statuses listed above.

\section{Activation Failure and Recovery}

A package can be valid and persistently loaded without enough RAM being
available to activate it. Loading its bytes does not reserve its future
execution RAM. If subsequent installation cannot allocate that RAM, it returns
\texttt{6A84} and does not create an application instance; the loaded package
remains available. The kernel can then process another selection or installation
request. An allocation failure when reactivating an existing instance does
not, by itself, delete its persistent registry record.

This distinction concerns resources allocated by the kernel for activation.
An allocation failure inside an already running Rustlet may instead be handled
by the application or terminate its execution through the runtime error path;
it is not automatically translated into \texttt{6A84}.

The \texttt{dyn\_rustlet} campaign exercises an executable whose declared stack
cannot fit on the test target, expects \texttt{6A84} on installation, checks
that no instance was created, and executes another valid Rustlet afterwards.
The same campaign separately expects \texttt{6F01} for isolated execution
faults. Both paths are exercised with the devkit and SCP03 management profiles.

\clearpage
\section{Implementation Sources}

The inventory is derived from these source groups:
\begin{itemize}
\item Transport and command dispatch:\newline
\path{kernel/firmware/src/apdu_manager.rs} and\newline
\path{kernel/firmware/src/main.rs}.
\item Security and activation: the secure-channel and Security Domain modules,
\path{kernel/firmware/src/selected_app.rs}, and\newline
\path{kernel/firmware/src/fae_runtime.rs}.
\item Built-in diagnostic modules:\newline
\path{kernel/firmware/src/kernel_main_app/}.
\item Runtime status definitions:\newline
\path{rustlets/rustlet_runtime/src/abi.rs}.
\end{itemize}
Changes to these response paths must update this appendix together with their
regression expectations.
